\section{Research Agenda}
\label{sec:research-agenda}

Revision 0.6 completes the first planned evidence stage and converts the nearest
remaining work into a public roadmap.

\subsection{Completed baseline: semantic mapping and fixed-input tests}

The ordered SWAP semantics are now documented, all $(v,b,e,r_E)$ branches are
enumerated, the fifth-wire reduced states are calculated, and the victim
subsystem is compared with a clean routed reference. The resulting theorem is
\cref{eq:caseii-fixed-input-theorem}. Information metrics and deterministic
regression outputs are included in the public artifact.

\subsection{Immediate priority: exact detector reconstruction}

The strongest next scientific task is to identify the full educational
circuit's actual acceptance register. This requires:
\begin{enumerate}
  \item naming every classical selector, measurement result, and comparison bit;
  \item reconstructing the detector flag as a Boolean function;
  \item producing its truth table over every branch;
  \item replacing the provisional rule in
  \cref{eq:caseii-provisional-acceptance}; and
  \item regenerating the deterministic outputs and manuscript claims.
\end{enumerate}
This step should precede any strong statement that the original educational
detector has been bypassed.

\subsection{Independent backend parity and software decomposition}

The NumPy implementation is the validated reference. A Qiskit implementation
should reproduce branch probabilities, reduced states, wire order, and metrics
within documented tolerances. In parallel, the legacy monolithic analysis file
should be decomposed into independently testable circuit, measurement, metric,
validation, plotting, and orchestration modules. These are engineering tasks,
but they strengthen the credibility and maintainability of the scientific
artifact.

\subsection{Formal information-flow property}

After the detector semantics are fixed, define a noninterference-style property
for reclaimed qubits. For a protected variable $X$, attacker register $E$, and
approved victim observation policy $\mathcal{V}$, a secure transformation should
establish both victim equivalence and an explicit information bound such as
\begin{equation}
  \mathcal{V}(\rho_{\mathrm{attack}})
  \approx \mathcal{V}(\rho_{\mathrm{baseline}}),
  \qquad I(X:E)\leq\varepsilon.
\end{equation}
The present case study supplies a counterexample to histogram-only verification,
not yet a complete compiler-level noninterference framework.

\subsection{Compiler prototype and hardware experiments}

A controlled malicious-reuse prototype should then perform ordinary liveness
analysis, identify reclaimed physical qubits, insert a benign synthetic payload,
and evaluate stronger contracts and trap circuits. Hardware and noise studies
should compare coherent cleanup, measurement-reset reuse, and fresh allocation,
including routing, crosstalk, leakage, and calibration drift. All experiments
should use synthetic secrets and authorized simulation or hardware.

\subsection{Venue and audience}

The work sits most naturally at the intersection of quantum software and quantum
security, with an educational circuit serving as the motivating case study. A
workshop, short research paper, or position paper aimed at that intersection is
a plausible first venue. Final venue selection should follow detector
reconstruction and external technical review.
